% Le terme d'alignement : deux pixels voisins, chacun avec sa direction propre e_1.
% À gauche les directions concordent, l'arête est forte ; à droite elles se croisent,
% l'arête s'effondre. Rouge Inria = ce qui est retenu, gris = ce qui est écarté.
    \begin{tikzpicture}[x=1cm,y=1cm,scale=1.0,every node/.style={scale=1.0},
        crete/.style={line width=7pt, gris_fonce_inria!45, line cap=round},
        dirv/.style={-{Latex[length=1.9mm]}, line width=1pt},
        pix/.style={circle,fill=inria-2024-bleu-canard,inner sep=0pt,minimum size=3.8pt}]

      % ---------------- directions concordantes ----------------
      \begin{scope}
        \draw[crete] (-0.15,-0.30)--(2.75,0.62);
        \node[pix] (i1) at (0.72,-0.02) {};
        \node[pix] (j1) at (1.88,0.35) {};
        \draw[dirv, inria-rouge] (0.72,-0.02)--++(0.62,0.20);
        \draw[dirv, inria-rouge] (1.88,0.35)--++(0.62,0.20);
        \draw[inria-rouge, line width=1.6pt] (i1)--(j1);
        \node[font=\scriptsize] at (0.60,-0.36) {$\theta(x_i)$};
        \node[font=\scriptsize] at (2.16,0.02) {$\theta(x_j)$};
        \node[font=\footnotesize, align=center] at (1.30,1.12)
              {$\theta(x_i)\simeq\theta(x_j)$\\[-1pt] {\scriptsize arête \textbf{forte}}};
      \end{scope}

      % ---------------- directions qui se croisent ----------------
      \begin{scope}[shift={(4.35,0)}]
        \draw[crete] (-0.15,-0.30)--(2.00,0.40);
        \draw[crete] (1.30,0.92)--(2.32,-0.48);
        \node[pix] (i2) at (0.72,-0.02) {};
        \node[pix] (j2) at (1.66,0.43) {};
        \draw[dirv, inria-rouge] (0.72,-0.02)--++(0.59,0.19);
        \draw[dirv, inria-rouge] (1.66,0.43)--++(0.36,-0.50);
        \draw[gris_fonce_inria!45, line width=0.8pt] (i2)--(j2);
        \node[font=\scriptsize] at (0.60,-0.36) {$\theta(x_i)$};
        \node[font=\scriptsize, anchor=west] at (2.10,-0.14) {$\theta(x_j)$};
        \node[font=\footnotesize, align=center] at (1.20,1.12)
              {$\theta(x_i)\neq\theta(x_j)$\\[-1pt] {\scriptsize arête \textbf{effondrée}}};
      \end{scope}
    \end{tikzpicture}
